% Generated draft: requires mathematical and source review.
\documentclass[11pt]{article}
\usepackage[margin=25mm]{geometry}
\usepackage{fontspec}
\setmainfont{Latin Modern Roman}
\usepackage{amsmath,amssymb,mathtools,unicode-math}
\setmathfont{Latin Modern Math}
\usepackage{enumitem,needspace,xurl,hyperref}
\setlength{\parindent}{0pt}
\setlength{\parskip}{5pt}
\setlength{\emergencystretch}{3em}
\newcommand{\N}{\mathbb{N}}
\newcommand{\Z}{\mathbb{Z}}
\newcommand{\Q}{\mathbb{Q}}
\newcommand{\R}{\mathbb{R}}
\newcommand{\C}{\mathbb{C}}
\newcommand{\F}{\mathbb{F}}
\newcommand{\A}{\mathbb{A}}
\newcommand{\PP}{\mathbb{P}}
\newcommand{\E}{\mathbb{E}}
\newcommand{\eps}{\varepsilon}
\newcommand{\dd}{\,\mathrm{d}}
\begin{document}
\section*{UnsolvedMath ID: 30004258}
\textbf{Source problem code:} catalogue entry
\subsection{Definitions and mathematical statement}
Let \(H\) be a reductive group over a field \(F\), acting on a space \(V\). The source denotes its nilpotent cone by \(\mathcal N_V\) and a Moy--Prasad neighborhood by \(V_{r+}\). These are source designations: the record does not specify the field hypotheses, the nature of \(V\) or its action, or the constructions of the cone and neighborhoods.

The proposed depth-controlled assertion is the following: for every instance of the intended setting and every \(H\)-distinguished representation \(\pi\) of depth \(r\), the relative local character expansion of \(\pi\) is valid on \(V_{r+}\). The ambient group represented by \(\pi\), the meaning of distinction, and the precise expansion are not specified in the record.

The source presents as an equivalent formulation the assertion that, for every \(r\in\R\),
\[
\operatorname{res}_{D(V)_{r+}}\widetilde J(V)_{r+}
=
\operatorname{res}_{D(V)_{r+}}J(\mathcal N_V).
\]
For a vector space \(E\), a subspace \(D\subseteq E\), and a collection \(\mathcal T\) of linear functionals on \(E\), restriction of the collection means
\[
\operatorname{res}_{D}\mathcal T
:=\{T|_D:T\in\mathcal T\}.
\]
This explains restriction notation only if the displayed source objects have these types. The record does not define \(D(V)_{r+}\), \(\widetilde J(V)_{r+}\), or \(J(\mathcal N_V)\), so their intended domains and membership conditions remain unresolved. In particular, the asserted equivalence and the relation between the quantifier over all real \(r\) and representation depths require clarification. The displayed equality is preserved as the source's proposed target, not asserted as a theorem. A fully specified mathematical conjecture requires the missing hypotheses and definitions.
\subsection{Short English statement}
The problem asks whether the depth of an \(H\)-distinguished representation controls a neighborhood on which its relative local character expansion holds. It also proposes an equivalent equality of restricted distribution spaces for every real parameter. The supplied record leaves the ambient representation setting and several essential spaces undefined, so this remains a statement-only draft requiring mathematical review.
\subsection{Sources}
\begin{itemize}
\item Local UnsolvedMath registry, supplied source record: permanent ID \(30004258\), title ``Depth-Controlled Relative Local Character Expansions,'' source problem code \texttt{catalogue entry}, in the Oberwolfach Reports Open Problems collection. The record gives proposed year 2019 and names no individual proposer.
\item Supplied report citation: \emph{New Developments in Representation Theory of p-adic Groups} (2020), \emph{Oberwolfach Reports}, pp.\ 2739--2819, DOI \texttt{10.4171/owr/2019/44}. Supplied source link: \url{https://doi.org/10.4171/owr/2019/44}. The background describes this collection as extracted from workshop reports and their open-problem discussions; no problem-specific page is supplied.
\item Supplied literature reference: ``OWR 2019 44,'' \url{https://ems.press/content/serial-article-files/46824}.
\item Literature assessment dated \texttt{2026-08-22}: the supplied curation labels the problem partially solved and states that relative local character expansion is known on some nilpotent neighborhoods and in favorable symmetric or Galois cases, while the asserted depth-controlled neighborhood statement is not proved in general. This is supplied dated curation derived from the OWR record, not an independently verified literature assessment; its time-sensitive status and cited claims require independent verification.
\item The registry labels the statement \texttt{corrected\_verified}, explaining that mathematical content was preserved while notation, typography, and sentence structure were repaired for publication. This is supplied editorial metadata, not independent mathematical verification of the formulation or its claimed equivalence.
\end{itemize}
\end{document}
